\documentclass[11pt,a4paper]{article}
\usepackage[T1]{fontenc}
\usepackage[latin1]{inputenc}
\usepackage{lmodern}
\usepackage{float}
\usepackage{a4wide}
\usepackage[dvips]{graphicx}

\usepackage[
pdfauthor={ACE Project Team},
pdftitle={Developers Guide},
pdfcreator={pdftex},
]{hyperref}

\usepackage{sectsty}
\allsectionsfont{\sffamily}

\usepackage{fancyheadings} 
\pagestyle{fancy} 
\lhead{\textsf{\textbf{ACE} \\ \small{a collaborative editor}}}
\chead{}
\rhead{
\parbox[c]{3cm}{\includegraphics[height=0.875cm,width=3cm]{../../images/logo_BFH.eps}}
\parbox[c]{2.2cm}
{\tiny{\textsf{Berner Fachhochschule \\
Hochschule f�r \\
Technik und Informatik}}}}
\lfoot{}
\cfoot{\textsf{\thepage}}
\rfoot{}
\setlength{\headrulewidth}{0.6pt}
\setlength{\footrulewidth}{0.6pt}
\setlength{\topmargin}{-50pt}
\addtolength{\headheight}{50pt}

\usepackage{colortbl}

\newcommand{\headercol}[2]{\multicolumn{1}{|>{\bfseries\columncolor[gray]{0.82}}p{#1}|}{\textsf{#2}}}
\newcommand{\ace}[0]{\emph{ACE }}



\begin{document}

\begin{titlepage}
\thispagestyle{empty}
  \includegraphics[height=1.5in]{../images/pix.eps}

  \begin{center}

    {\fontsize{40}{45} \textbf{\textsf{ACE}}} \\
    \textsf{a collaborative editor} \\
        
    \vspace{36pt}
        
    {\huge{\textbf{\textsf{}}}} \\

    \vspace{36pt}

	\textsf{Berne University of Applied Sciences} \\
    \textsf{School of Engineering and Information Technology} \\
    
  \end{center}

  \vfill
  
  \begin{tabular}{ll}
   \hline

   \\

   \multicolumn{1}{>{\bfseries}p{1.5in}}{\textsf{Date:}} &
   \multicolumn{1}{>{}p{4.3in}}{\textsf{08.11.2005}}          \\
   
   \\
   
   \multicolumn{1}{>{\bfseries}p{1.5in}}{\textsf{Version:}}     &   
   \multicolumn{1}{>{}p{4.3in}}{\textsf{0.1}}                 \\

   \\
   
   \multicolumn{1}{>{\bfseries}p{1.5in}}{\textsf{Projectteam:}}                 &
   \multicolumn{1}{>{}p{4.3in}}{\textsf{Mark Bigler (biglm2@hta-bi.bfh.ch)}}  \\
   \multicolumn{1}{>{\bfseries}p{1.5in}}{}                                      &
   \multicolumn{1}{>{}p{4.3in}}{\textsf{Simon Raess (rasss@hta-bi.bfh.ch)}}    \\
   \multicolumn{1}{>{\bfseries}p{1.5in}}{}                                      &
   \multicolumn{1}{>{}p{4.3in}}{\textsf{Lukas Zbinden (zbinl@hta-bi.bfh.ch)}} \\   
   
   \\
   
   \multicolumn{1}{>{\bfseries}p{1.5in}}{\textsf{Receivers:}}                       &
   \multicolumn{1}{>{}p{4.3in}}{\textsf{Jean-Paul Dubois (doj@hta-bi.bfh.ch)}}       \\
   \multicolumn{1}{>{\bfseries}p{1.5in}}{}                                          &
   \multicolumn{1}{>{}p{4.3in}}{\textsf{Claude Fuhrer (frc@hta-bi.bfh.ch)}}       \\

   \\
   
   \multicolumn{1}{>{\bfseries}p{1.5in}}{\textsf{Location:}}               &   
   \multicolumn{1}{>{}p{4.3in}}{\textsf{Subversion Repository}} \\

   \\  
   
   \hline
  \end{tabular}

\end{titlepage}


\tableofcontents


\section{Introduction}

The \emph{Developers Guide} contains information on how to build the application
from source, how to install, and run the application. Further, information
about the subversion repository is given.

Unless you want to change the source code, we do not recommend to build the
application from source. We provide installers for Windows and a double
clickable application for Mac OS X. These install a completely working
application without the need for all the details that follow.


\section{Setup}
ACE uses Apache Ant (see \href{http://ant.apache.org/}{http://ant.apache.org/})
to build the application. If you have not installed Ant, download and install
it from the above URL. Further, we use a special Ant task to download the
libraries used by ACE. These libraries are not stored inside of the source
repository.

\subsection{Maven 2.0 Ant Tasks}
Maven 2.0 now comes with a set of Ant tasks that can be used to utilise Maven's
artifact handling features from within Ant. This includes most notably the
transitive \emph{Dependency management} handling of Maven 2.0. You can
download these Ant tasks from 
\href{http://maven.apache.org/download.html}{http://maven.apache.org/download.html}.

\subsubsection{Installation}
To install these Ant tasks you have to do the following:
\begin{itemize}
 \item download the maven-artifact-ant-2.0-dep.jar from the Maven download site
 \item copy the downloaded jar file to the directory \texttt{ANT_HOME/lib} where \texttt{ANT_HOME} is where Ant is installed
\end{itemize}

\subsubsection{Maven Repositories}
ACE uses two different Maven repository. A Maven repository is a place where
dependencies are installed. The first repository is a custom repository that
stores all dependencies that are not available from the default Maven
repository. The URL of this repository is http://ace.iserver.ch/maven2. The
second repository is the standard Maven repository at the URL
http://repo1.maven.org/maven2/.

\subsubsection{Dependencies}
In the \texttt{build.xml} you find the declaration of the dependencies.

\small{\begin{verbatim}
<artifact:dependencies pathId="dependency.classpath">
		<dependency groupId="jdom" 
		            artifactId="jdom" 
		            version="1.0"/>
		<dependency groupId="commons-beanutils" 
		            artifactId="commons-beanutils" 
		            version="1.7.0"/>
        ...
</artifact:dependencies>
\end{verbatim}}

A dependency has a group id, an artifact id, and a version. You have to know
these in order to add dependencies (browse the Maven repositories to find
out the correct values). The dependencies are downloaded to a local repository,
which can be found in the directory \texttt{~/.m2/repository/}.

The id given to the set of dependencies with the \texttt{pathId} attribute
can be used like any other path in the Ant build file. The Maven Ant tasks
allow to get the correct versions of dependencies without needing to place
them in the source repository.

To find out more about the Maven Ant tasks visit
\href{http://maven.apache.org/ant-tasks.html}{http://maven.apache.org/ant-tasks.html}.

\subsubsection{First Time Build}
When you run the build the first time you will get a couple of messages like:

\small{\begin{verbatim}
[artifact:dependencies] Downloading: log4j/log4j/1.2.12/log4j-1.2.12.pom
[artifact:dependencies] Transferring 0K
[artifact:dependencies] Downloading: log4j/log4j/1.2.12/log4j-1.2.12.jar
[artifact:dependencies] Transferring 349K
\end{varbatim}}

You will get also some warnings, because the task first checks the custom
repository of ACE, which does not contain all the dependencies of the project.


\section{Directory Structure}

ACE has the following contents in the top-level directory:

\begin{table}[H]
 \begin{tabular}{|l|l|}
  \hline
  /.classpath         &  Eclipse classpath file    \\
  \hline
  /.project           &  Eclipse project file      \\
  \hline
  /build              &  build products            \\
  \hline
  /build.xml          &  Ant build file            \\
  \hline
  /dist               &  distribution directory    \\
  \hline
  /doc                &  documentation             \\
  \hline
  /project.properties &  Maven project properties  \\
  \hline
  /project.xml        &  Maven project model       \\
  \hline
  /src                &  source files              \\
  \hline
  /www                &  the project website       \\
  \hline
 \end{tabular}
 \caption{Top-level directory structure}
\end{table}

\subsection{Build Products}
The \texttt{build} folder contains the build results of the project:

\begin{table}[H]
 \begin{tabular}{|l|l|}
  \hline
  /build/api             &  generated Javadoc API  \\
  \hline
  /build/classes         &  class files            \\
  \hline
  /build/resources       &  resource files         \\
  \hline
  /build/testreports     &  JUnit test reports     \\
  \hline
 \end{tabular}
 \caption{build directory structure}
\end{table}

\subsection{Sources}
The \texttt{src} folder contains the sources of the project:

\begin{table}[H]
 \begin{tabular}{|l|l|}
  \hline
  /src/installer
  &  NSIS Windows installer        \\
  \hline
  /src/integration-test
  &  integration tests     \\
  \hline
  /src/java
  &  Java source files \\
  \hline
  /src/resources
  &  resources used by the application \\
  \hline
  /src/stubs 
  &  stub classes used for testing \\
  \hline
  /src/test
  &  JUnit tests \\
  \hline
  /src/test-app
  &  simple test applications to quickly test parts of the application \\
  \hline
 \end{tabular}
 \caption{source directory structure}
\end{table}

\subsection{Documentation}
The \texttt{doc} folder contains the sources of the project:

\begin{table}[H]
 \begin{tabular}{|l|l|}
  \hline
  /doc/developersguide
  &  the \LaTeX{} source for the developers guide \\
  \hline
  /doc/finalreport
  &  the \LaTeX{} source for the final report \\
  \hline
  /doc/images
  &  images used in the documentation \\
  \hline
  /doc/latex
  &  the \LaTeX{} templates \\
  \hline
  /doc/projectmanual
  &  the \LaTeX{} source for the project manual \\
  \hline
  /doc/systemrequirements
  &  the \LaTeX{} source for the system requirements \\
  \hline
  /doc/templates
  &  source code templates \\
  \hline
  /doc/usermanual
  &  the \LaTeX{} source for the system user manual \\
  \hline
 \end{tabular}
 \caption{documentation directory structure}
\end{table}


\section{Build File}

The Ant build file contains the all the relevant targets. The main targets
are:

\begin{itemize}
 \item \texttt{compile}
 \item \texttt{test}
 \item \texttt{integration-test}
 \item \texttt{run-gui}
 \item \texttt{clean}
 \item \texttt{javadoc}
 \item \texttt{dist}
\end{itemize}

\subsection{Compiling the Sources}
To compile the sources you use the Ant target \texttt{compile}. 

\subsection{Running the Tests}
There are two test targets. The first that runs the unit tests is called
\texttt{test}. These tests are run regularly by our continous integration
server. The test sources are found in the directory \texttt{/src/test}. They
must run quickly and deterministically.

The other test target is \texttt{integration-test} which runs integration tests.
Integration tests test large parts of the system and thus run slower than
the pure unit tests.

\subsection{Running the Application}
You can run the application by calling the \texttt{run-gui} task. This starts
the main class of the whole application. If you want to run another class
with the main argument, you can start it with by calling the \texttt{run}
task with a system property \texttt{class} set to the correct class name.
You do this from the command line with an additional argument 
\texttt{-Dclass=...}.

\subsection{Cleaning the Build Results}
The well-known \texttt{clean} target can be used to clean the build directory,
forcing a complete rebuild.

\subsection{Javadoc}
The \texttt{javadoc} target creates the Javadoc of the application into the
directory \texttt{/build/api}.

\subsection{Creating a Distribution}
The Ant target \texttt{dist} creates a \texttt{tar.gz} file containing the
sources, the javadoc, and all the necessary jar files into the directory
\texttt{/dist}.


\section{Building the Installers}


\end{document}
